\begin{comment}
\begin{table*}[!htb]
    \centering
    \begin{threeparttable}
    \begin{tabular}{|p{45pt}|c|c|c|c|c|c|c|c|c|c|}
    \hline
        \multirow{3}*{\raisebox{-0.5\height}{\textbf{Method}}} & \multirow{3}*{\raisebox{-0.5\height}{\textbf{LLM}}} & \textbf{NQ-} & \multicolumn{2}{c|} {\textbf{Trivia}} & \multicolumn{2}{c|}{\textbf{Web}} & \multicolumn{2}{c|}{\textbf{Hotpot}} & \multicolumn{2}{c|}{\textbf{2WikiMulti}} \\
        && \textbf{open} & \multicolumn{2}{c|} {\textbf{QA}} & \multicolumn{2}{c|}{\textbf{Questions}} & \multicolumn{2}{c|}{\textbf{QA}} & \multicolumn{2}{c|}{\centering{\textbf{HopQA}}} \\
        %& & \centering{\textbf{NQ-open}} & \multicolumn{2}{c|}{\raisebox{-0.5\height}{\textbf{TriviaQA}}} & \multicolumn{2}{c|}{\raisebox{-0.5\height}{\textbf{WebQuestions}}} & \multicolumn{2}{c|}{\raisebox{-0.5\height}{\textbf{HotpotQA}}} & \multicolumn{2}{p{50pt}|}{\centering{\textbf{2WikiMulti HopQA}}} \\
        \cline{3-11}
        && EM & EM & F1 & EM & F1 & EM & F1 & EM & F1 \\
        \hline
        \multirow{3}{40pt}{No-Retrieval LLM} & GPT & 35.06 & 44.38 & 48.15 & 37.43 & 45.81 & 28.41 & 34.68 & 41.63 & 49.13 \\
        & Q & 33.31 & 42.16 & 45.74 & 35.56 & 43.52 & 25.57 & 31.21 & 37.47 & 44.22 \\
        & DS & 36.11 & 45.71 & 49.59 & 38.55 & 47.18 & 29.83 & 36.41 & 43.71 & 51.59 \\
        \hline
        \multirow{3}{40pt}{Vanilla RAG} & GPT & 48.35 & 58.31 & 63.52 & 44.52 & 57.63 & 39.14 & 44.85 & 52.33 & 57.71 \\
        & Q & 39.78 & 55.39 & 60.34 & 42.29 & 54.75 & 35.23 & 40.36 & 47.1 & 51.94 \\
        & DS & 43.13 & 60.06 & 65.43 & 45.86 & 59.36 & 41.1 & 47.09 & 54.95 & 60.6 \\
        \hline
        \multirow{3}{40pt}{CoT-on-Concat} & GPT & 49.25 & 60.57 & 67.38 & 44.91 & 56.42 & 39.36 & 45.38 & 52.55 & 58.24\\
        & Q & 41.96 & 57.54 & 64.01 & 42.66 & 53.6 & 35.42 & 40.84 & 47.3 & 52.42\\
        & DS & 45.5 & 62.39 & 69.4 & 46.26 & 58.11 & 41.33 & 47.65 & 55.18 & 61.15\\
        \hline
        \multirow{3}{40pt}{Wikipedia Graph} & GPT & 49.68 & 54.81 & 59.71 & 42.37 & 55.88 & 51.76 & 60.14 & 55.52 & 64.51\\
        & Q & 42.17 & 52.07 & 56.72 & 40.25 & 53.09 & 46.58 & 54.13 & 49.97 & 58.06\\
        & DS & 45.72 & 56.45 & 61.5 & 43.64 & 57.56 & 54.35 & 63.15 & 58.3 & 67.74\\
        \hline
        \multirow{3}{40pt}{FiD} & GPT & 53.62 & 72.61 & 76.51 & 49.30 & 62.81 & 56.10 & 62.51 & 56.14 & 65.62\\
        & Q & 50.94 & 68.98 & 72.68 & 46.83 & 59.67 & 50.49 & 56.26 & 50.53 & 59.06\\
        & DS & 55.23 & 74.79 & 78.81 & 50.78 & 64.69 & 58.9 & 65.64 & 58.95 & 68.9\\
        \hline
        \multirow{3}{40pt}{ReAct} & GPT & 51.45 & 63.97 & 74.24 & 42.91 & 55.14 & 56.27 & 64.20 & 56.52 & 66.52\\
        & Q & 47.51 & 60.77 & 70.53 & 40.76 & 52.38 & 50.64 & 57.78 & 50.87 & 59.87\\
        & DS & 51.51 & 65.89 & 76.47 & 44.2 & 56.79 & 59.08 & 67.41 & 59.35 & 69.85\\
        \hline
        \multirow{3}{40pt}{IRCoT} & GPT & 52.3 & 69.21 & 75.04 & 46.36 & 59.62 & 44.51 & 55.14 & 57.07 & 68.32\\
        & Q & 49.68 & 65.75 & 71.29 & 44.04 & 56.64 & 40.06 & 49.63 & 51.36 & 61.49\\
        & DS & 53.87 & 71.29 & 77.29 & 47.75 & 61.41 & 46.74 & 57.9 & 59.92 & 71.74\\
        \hline
        \multirow{3}{40pt}{SELF-RAG} & GPT & 53.14 & 70.38 & 78.88 & 48.96 & 61.58 & 43.61 & 51.79 & 58.71 & 69.50\\
        & Q & 50.48 & 66.86 & 74.94 & 46.51 & 58.5 & 39.25 & 46.61 & 52.84 & 62.55\\
        & DS & 54.73 & 72.49 & 81.25 & 50.43 & 63.43 & 45.79 & 54.38 & 61.65 & 72.98\\
        \hline
        \multirow{3}{40pt}{Multi-Hop RAG} & GPT & 50.15 & 68.76 & 74.81 & 43.12 & 56.97 & 54.17 & 63.42 & 56.91 & 66.35\\
        & Q & 45.89 & 65.32 & 71.07 & 40.96 & 54.12 & 48.75 & 57.08 & 51.22 & 59.72\\
        & DS & 49.76 & 70.82 & 77.05 & 44.41 & 58.68 & 56.88 & 66.59 & 59.76 & 69.67\\
        \hline
        \multirow{3}{40pt}{\textbf{Ours}} & GPT & \textbf{53.85} & \textbf{71.12} & \textbf{79.45} & \textbf{50.15} & \textbf{63.24} & \textbf{57.42} & \textbf{66.85} & \textbf{59.85} & \textbf{70.65}\\
        & Q & \textbf{51.25} & \textbf{67.95} & \textbf{75.82} & \textbf{47.65} & \textbf{59.85} & \textbf{52.15} & \textbf{59.45} & \textbf{53.45} & \textbf{63.85}\\
        & DS & \textbf{55.65} & \textbf{73.15} & \textbf{81.85} & \textbf{51.55} & \textbf{65.15} & \textbf{60.45} & \textbf{68.95} & \textbf{62.85} & \textbf{74.15}\\
        \hline
    \end{tabular}
    \begin{tablenotes}
        \small
        \item *\textit{GPT}, \textit{Q}, and \textit{DS} refer to GPT-4o-mini, Qwen-turbo, and DeepSeek-V3, respectively. 
    \end{tablenotes}
    \caption{RAG evaluation on question-answering benchmark dataset.}
    \label{tab:main1}
    \end{threeparttable}
\end{table*}
\end{comment}


This section elaborates on the evaluation of our SLM-RAG method.
We utilize Llama-3.2-1B-Instruct as SLM.
This small-sized language model is sufficient to function as a semantic controller for all related subtasks in our method.
This benefit arises more from the instruction-following consistency than the model scale.
Unless otherwise specified, we set retrieval depths at $m=20$ and $k=10$.
Additionally, we set completeness threshold at $\delta=0.5$ and coverage threshold at $\epsilon=0.6$ (further details in Appendix~\ref{sec:threshold}).


\begin{table*}[!htb]
\small
    \centering
    \begin{threeparttable}
    \begin{tabular}{|p{50pt}|c|c|c|c|c|c|c|}
    \hline
        \multirow{2}*{\raisebox{-0.5\height}{\textbf{Method}}} & \multirow{2}*{\raisebox{-0.5\height}{\textbf{LLM}}} & \multicolumn{3}{c|} {\textbf{ELI5}} & \multicolumn{3}{c|}{\textbf{FEVER}} \\
        %& & \centering{\textbf{NQ-open}} & \multicolumn{2}{c|}{\raisebox{-0.5\height}{\textbf{TriviaQA}}} & \multicolumn{2}{c|}{\raisebox{-0.5\height}{\textbf{WebQuestions}}} & \multicolumn{2}{c|}{\raisebox{-0.5\height}{\textbf{HotpotQA}}} & \multicolumn{2}{p{50pt}|}{\centering{\textbf{2WikiMulti HopQA}}} \\
        \cline{3-8}
        && ROUGE-1 & ROUGE-2 & ROUGE-L & Acc & FEVER Score & F1 \\
        \hline
        \multirow{3}{50pt}{w/o RAG} & GPT & 23.62 & 4.25 & 20.51 & 61.97 & 50.87 & 64.52\\
        & Qwen & 21.73 & 3.91 & 18.87 & 58.87 & 48.33 & 61.29 \\
        & DeepSeek & 23.15 & 4.16 & 20.1 & 63.83 & 52.4 & 66.46 \\
        \hline
        \multirow{3}{40pt}{Na\"ive RAG} & GPT & 30.21 & 6.77 & 24.32 & 66.82 & 62.39 & 77.17 \\
        & Qwen & 27.79 & 6.23 & 22.37 & 63.48 & 59.27 & 73.31 \\
        & DeepSeek & 29.61 & 6.63 & 23.83 & 68.82 & 64.26 & 79.49 \\
        \hline
        \multirow{3}{40pt}{CoT-on-Concat} & GPT & 31.20 & 8.39 & 25.31 & 68.91 & 66.34 & 78.35 \\
        & Qwen & 28.7 & 7.72 & 23.29 & 65.46 & 63.02 & 74.43 \\
        & DeepSeek & 30.58 & 8.22 & 24.8 & 70.98 & 68.33 & 80.7 \\
        \hline
        \multirow{3}{40pt}{Wikipedia Graph} & GPT & 32.01 & 9.15 & 26.91 & 74.39 & 67.31 & 80.76\\
        & Qwen & 29.45 & 8.42 & 24.76 & 70.67 & 63.94 & 76.72\\
        & DeepSeek & 31.37 & 8.97 & 26.37 & 76.62 & 69.33 & 83.18\\
        \hline
        \multirow{3}{40pt}{FiD} & GPT & 34.72 & 10.52 & 27.37 & 80.09 & 71.39 & 86.31\\
        & Qwen & 31.94 & 9.68 & 25.18 & 76.09 & 67.82 & 81.99\\
        & DeepSeek & 34.03 & 10.31 & 26.82 & 82.49 & 73.53 & 88.9\\
        \hline
        \multirow{3}{40pt}{ReAct} & GPT & 30.21 & 7.26 & 24.81 & 77.52 & 69.37 & 83.03\\
        & Qwen & 27.79 & 6.68 & 22.83 & 73.64 & 65.9 & 78.88\\
        & DeepSeek & 29.61 & 7.11 & 24.31 & 79.85 & 71.45 & 85.52\\
        \hline
        \multirow{3}{40pt}{IRCoT} & GPT & 31.28 & 9.63 & 27.34 & 82.61 & 74.28 & 87.01\\
        & Qwen & 28.78 & 8.86 & 25.15 & 78.48 & 70.57 & 82.66\\
        & DeepSeek & 30.65 & 9.44 & 26.79 & 85.09 & 76.51 & 89.62\\
        \hline
        \multirow{3}{40pt}{SELF-RAG} & GPT & 34.09 & 12.83 & 30.20 & 83.62 & 79.68 & 89.92\\
        & Qwen & 31.36 & 11.8 & 27.78 & 79.44 & 75.7 & 85.42\\
        & DeepSeek & 33.41 & 12.57 & 29.6 & 86.13 & 82.07 & 92.62\\
        \hline
        \multirow{3}{40pt}{Multi-Hop RAG} & GPT & 27.86 & 6.32 & 23.61 & 73.62 & 68.67 & 82.00\\
        & Qwen & 25.63 & 5.81 & 21.72 & 69.94 & 65.24 & 77.9\\
        & DeepSeek & 27.3 & 6.19 & 23.14 & 75.83 & 70.73 & 84.46\\
        \hline
        \multirow{3}{40pt}{\textbf{InSemRAG (Ours)}} & GPT & \textbf{35.42} & \textbf{13.45} & \textbf{31.15} & \textbf{85.12} & \textbf{81.25} & \textbf{91.05}\\
        & Qwen & \textbf{32.85} & \textbf{12.45} & \textbf{28.95} & \textbf{81.15} & \textbf{77.25} & \textbf{86.85}\\
        & DeepSeek & \textbf{34.85} & \textbf{13.15} & \textbf{30.85} & \textbf{87.45} & \textbf{83.65} & \textbf{93.55}\\
        \hline
    \end{tabular}
    %\begin{tablenotes}
    %    \small
    %    \item *\textit{GPT}, \textit{Q}, and \textit{DS} refer to GPT-4o-mini, Qwen-turbo, and DeepSeek-V3, respectively. 
    %\end{tablenotes}
    \caption{RAG evaluation on long-form generation and fact verification benchmarks. Our method yields competitive performance over strong reasoning-centric baselines.}
    \label{tab:main2}
    \end{threeparttable}
\end{table*}


\subsection{Main Results}
We evaluate the RAG performance on three LLMs as generators.
Specifically, we utilize GPT-4o-mini, Qwen-turbo, and DeepSeek-V3 to cover diverse architectures and deployment settings, enabling a fair evaluation of robustness and cross-model generalization.
%NQ-Open is an open-domain natural QA benchmark of 1k test set, TriviaQA comprises 17k complex and compositional QA sets, and WebQuestions is a QA benchmark for evaluating structured knowledge bases with 2k test set.
We compare our method against several strong baselines such as generation without RAG, Na\"ive RAG~\citep{lewis2020retrieval, gao2023retrieval}, CoT-on-Concat~\cite{wei2022chain}, Wikipedia Graph~\cite{lin2025structured}, FiD~\cite{izacard2021leveraging}, ReAct~\cite{yao2022react}, IRCoT~\cite{trivedi2023interleaving}, SELF-RAG~\cite{asai2024self}, and Multi-Hop RAG~\cite{tang2024multihop}.
We evaluate on five question-answering (QA) benchmark datasets (i.e, NQ-open~\cite{kwiatkowski2019natural}, TriviaQA~\cite{joshi2017triviaqa}, WebQuestions~\cite{berant2013semantic}, HotpotQA~\cite{yang2018hotpotqa}, and 2WikiMultiHopQA~\cite{ho2020constructing}), long-form generation benchmark (ELI5~\cite{fan2019eli5}), and fact verification benchmark (FEVER~\cite{thorne2018fever}).
We measure Exact Match (EM) and F1 scores for QA tasks, ROUGE scores~\cite{lin2004rouge} for long-form generation task, as well as accuracy, FEVER Score~\cite{thorne2018fever}, and F1 score on fact verification task.
Table~\ref{tab:main1} presents the results on QA tasks and Table~\ref{tab:main2} shows the evaluation on long-form generation and fact verification tasks.
We observe that across all generators, our method consistently outperforms Na\"ive RAG and CoT-on-Concat on both single- and multi-hop QA, demonstrating that dual-view query enrichment and dynamic sparse–dense routing improve evidence recall without retriever retraining. 
These benefits are most pronounced on entity-centric benchmarks (NQ-open, TriviaQA, WebQuestions), where lexical precision and semantic intent must be jointly balanced. 
Performance gains further amplify on multi-hop and evidence-sensitive tasks (HotpotQA, 2WikiMultiHopQA, ELI5, FEVER), where fixed-size chunking often causes semantic fragmentation. 
In these settings, integrity-aware evidence construction yields substantial improvements over strong reasoning-centric baselines (ReAct, IRCoT, SELF-RAG). 
Consistent gains on FEVER and long-form generation also indicate improved evidential faithfulness and global coherence. 
%Overall, adaptive routing improves what evidence is retrieved, while integrity-aware construction improves how it is structured and consumed, with both jointly driving the observed improvements.












\subsection{Ablation Studies}

\pb{Generalizability and Modular Contribution.}
To evaluate the generalizability of our framework, we swap our default SLM with Qwen2.5-1.5B-Instruct (SLM swap).
We also examine the contribution of each module to the RAG performance.
Specifically, we evaluate four modular conditioning: without query enrichment in IAR (w/o QE-IAR), without dynamic weighting (weighting retrieval channels with uniform weights $\alpha=\beta=\gamma$) in IAR (w/o DW-IAR), without candidate evidence refinement step in SPC (w/o ER-SPC), without SPC (passing the retrieved output of QAR module to LLM) (w/o SPC), and replacing SPC with heuristic rules (SPC $\xrightarrow{}$ heuristic).
Note that we utilize the following heuristic rules to determine the completeness of a chunk:
\begin{itemize}
    \item is terminated by terminal punctuation marks (i.e., ., ?, or !),
    \item contains balanced pairs of quotation marks, parentheses, and brackets, and
    \item has a minimum token length of 50.
\end{itemize}
We also test the effectiveness of our first retrieval depth value $m$ with respect to the final depth $k$ (depth $m=4k$).
Table~\ref{tab:genmod} demonstrates that there exists a significant performance gap between baseline and heuristic-based RAG (e.g., -5.60 on HotpotQA), demonstrating that semantic fragmentation often occurs behind grammatically "correct" endings, such as unresolved co-references (e.g., "He then decided...") or dangling logical premises, where SLM offers proper assistance in repairing.
%This result confirms that the gains arise from our RAG system itself rather than the choice of SLMs. 
Removing dynamic channel weighting in IAR consistently degrades performance, validating its role in robust evidence recall. 
Disabling query enhancement in IAR further reduces accuracy, particularly for multi-hop and abstract queries. 
The largest drops occur when SPC is removed, highlighting semantic fragmentation as a primary bottleneck in complex reasoning and long-form generation.
Bypassing candidate evidence refinement in SPC or replacing SPC with heuristic rules also significantly harms performance, confirming the necessity of semantic filtering and repair. 
Increasing $m$ yields diminishing returns, indicating that selection effectively balances recall and noise. 
Finally, similar results with alternative lightweight SLMs verify that the framework is model-agnostic. 
Overall, IAR governs what evidence is retrieved, while SPC governs how it is filtered and structured, and all components are jointly necessary for full performance gains.






\begin{table}[!htb]
\small
    \centering
    \begin{tabular}{|c|c|c|}
    \hline
        \multirow{2}*{\textbf{Evaluation}} & \textbf{HotpotQA} & \textbf{ELI5} \\
        & \textbf{(FI)} & \textbf{(ROUGE-L)} \\
        %\cline{3-8}
        \hline
        Baseline (Ours) & 66.85 & 31.15 \\
        \hline
        SLM swap & 65.42 & 30.22 \\
        w/o QE-IAR & 64.95 & 30.25 \\
        w/o DW-IAR & 64.1 & 29.8 \\
        w/o ER-SPC & 63.80 & 28.95 \\
        w/o SPC & 59.45 & 25.12 \\
        SPS $\xrightarrow{}$ heuristic & 61.25 & 27.8\\
        retrieval depth $m=4k$ & 66.20 & 31.05\\
        \hline
    \end{tabular}
    \caption{Evaluation on generalizability and modular contribution. Both IAR and SPC make substantial contributions to performance improvement.}
    \label{tab:genmod}
\end{table}



\pb{Robustness of Hyperparameter.}
We examine the model performance under varying chunk length $l$ and final retrieval depth $k$.
We evaluate on HotpotQA and measure the F1 score.
Table~\ref{tab:hyper} demonstrates that our method is persistently stronger than other RAG frameworks across chunk length and final retrieval depth.
This result demonstrates that semantic fragmentation severely degrades conventional RAG performance as chunk sizes shrink, while our method remains robust.
Na\"ive RAG and Multi-Hop RAG deteriorate sharply under small chunks, indicating that iterative reasoning cannot offset truncated evidence. 
In contrast, our approach consistently achieves the best F1 across all settings, with the performance gap widening as fragmentation increases. 
This demonstrates that integrity-aware evidence construction effectively repairs fragmented evidence and enables robust reasoning under adversarial retrieval conditions.


\begin{table}[!htb]
\small
    \centering
    \begin{tabular}{|m{40pt}|c|c|c|c|}
    \hline
        \multirow{2}*{\textbf{Method}} & \multirow{2}*{\textbf{$l$}} & \multicolumn{3}{c|}{\textbf{$k$}} \\
        \cline{3-5}
        && \textbf{5} & \textbf{10} & \textbf{15} \\
        \hline
        \multirow{4}{0pt}{Baseline (Ours)} & \textbf{512} & 58.52 & 66.85 & 68.45 \\
        & \textbf{256} & 56.8 & 65.15 & 67.05 \\
        & \textbf{128} & 51.45 & 60.85 & 63.4\\
        & \textbf{64} & 45.25 & 54.95 & 58.1\\
        \hline
        \multirow{4}{40pt}{Na\"ive} & \textbf{512} & 48.35 & 57.71 & 59.12 \\
        & \textbf{256} & 40.12 & 48.55 & 50.4 \\
        & \textbf{128} & 32.45 & 40.15 & 43.2\\
        & \textbf{64} & 24.15 & 31.8 & 34.55\\
        \hline
        \multirow{4}{40pt}{Multi-Hop RAG} & \textbf{512} & 54.17 & 63.42 & 65.2 \\
        & \textbf{256} & 50.25 & 60.12 & 62.45 \\
        & \textbf{128} & 45.6 & 55.45 & 58.12\\
        & \textbf{64} & 38.45 & 48.2 & 51.3\\
        \hline
    \end{tabular}
    \caption{Evaluation on robustness of hyperparameters. Our method is consistently competitive against other methods across various chunk length $l$ and final retrieval depth $k$.}
    \label{tab:hyper}
\end{table}



\pb{Latency Analysis.}
To evaluate the applicability of our framework to a real-case RAG process, we measure the latency of the end-to-end process on the HotpotQA benchmark.
For fair comparison, we set a fixed chunk length of 256.
Table~\ref{tab:latent} shows that our method operates at a competitive computation latency compared to other strong baselines.
Our method only slightly increases per-query latency over Na\"ive RAG while substantially outperforming it in F1, and remains far more efficient than Multi-Hop RAG, which is both slower and less accurate. 
This demonstrates that lightweight SLM-based coordination enables robust retrieval and evidence repair without costly multi-step reasoning.



\begin{table}[!htb]
\small
    \centering
    \begin{tabular}{|c|c|c|}
    \hline
        \textbf{Method} & \textbf{F1 Score} & \textbf{Latency (s/q)} \\
        %\cline{3-8}
        \hline
        Baseline (Ours) & 66.85 & 1.95 \\
        \hline
        Na\"ive RAG & 44.85 & 1.25 \\
        Multi-Hop RAG & 63.42 & 8.42 \\
        \hline
    \end{tabular}
    \caption{Evaluation on end-to-end latency. Our method offers competitive performance at 4$\times$ lower latency.}
    \label{tab:latent}
\end{table}


